\subsection{2.9 Diffusionsprojektion und der universelle Skalenparameter}

Nach der Auswahl der Diffusionsprojektion als vielversprechendstem Projektionsoperator begann die nächste Phase der numerischen Untersuchungen.

Die zentrale Frage lautete nun:

\[
\boxed{
\text{Ist das beobachtete } d_{\mathrm{eff}}\approx4
\text{ robust oder lediglich Zufall?}
}
\]

Insbesondere sollte untersucht werden, wie empfindlich die Projektion auf ihre freien Parameter reagiert.

%%%%%%%%%%%%%%%%%%%%%%%%%%%%%%%%%%%%%%%%%%%%%%%%%%%%%%%%%%%%%%

\section*{Der Diffusionsparameter}

Die Diffusionsprojektion basiert auf einem Diffusionsprozess auf dem relationalen Graphen.

Mathematisch kann dieser Prozess durch einen Operator

\[
P^t
\]

beschrieben werden.

Dabei bezeichnet

\[
t
\]

die effektive Diffusionszeit.

Physikalisch bedeutet dies:

\begin{itemize}

\item

kleines \(t\):
nur lokale Strukturen werden sichtbar,

\item

großes \(t\):
immer größere Bereiche verschmelzen zu globalen Strukturen.

\end{itemize}

Der Parameter

\[
t
\]

bestimmt somit unmittelbar die Beobachtungsskala.

%%%%%%%%%%%%%%%%%%%%%%%%%%%%%%%%%%%%%%%%%%%%%%%%%%%%%%%%%%%%%%

\section*{Der erste Diffusionsscan}

Im ersten systematischen Test wurde

\[
t
\]

über einen großen Bereich variiert.

Typische Ergebnisse waren:

\[
\begin{array}{c|c}
t
&
d_{\mathrm{eff}}
\\
\hline
0.25
&
39.8
\\
0.5
&
39.3
\\
1
&
36.8
\\
2
&
28.7
\\
4
&
13.5
\\
8
&
5.1
\end{array}
\]

Bereits dieses Resultat war außerordentlich aufschlussreich.

%%%%%%%%%%%%%%%%%%%%%%%%%%%%%%%%%%%%%%%%%%%%%%%%%%%%%%%%%%%%%%

\section*{Die wichtigste Beobachtung}

Die effektive Dimension nahm kontinuierlich mit wachsender Diffusionszeit ab.

Formal:

\[
t
\uparrow
\quad\Longrightarrow\quad
d_{\mathrm{eff}}
\downarrow.
\]

Dies bedeutet,

dass die Diffusionsprojektion tatsächlich Freiheitsgrade zusammenfasst.

Die Raumzeit entsteht also nicht plötzlich,

sondern entwickelt sich kontinuierlich aus dem hochdimensionalen OPP.

%%%%%%%%%%%%%%%%%%%%%%%%%%%%%%%%%%%%%%%%%%%%%%%%%%%%%%%%%%%%%%

\section*{Der Feinscan}

Da sich

\[
t
\approx
8
\]

bereits in der Nähe der gesuchten Vierdimension befand,

wurde anschließend ein deutlich feinerer Scan durchgeführt.

Untersucht wurden insbesondere

\[
t
=
6,7,8,9,10,11,12.
\]

Die Ergebnisse lauteten näherungsweise

\[
\begin{array}{c|c}
t
&
d_{\mathrm{eff}}
\\
\hline
6
&
7.79
\\
7
&
5.65
\\
8
&
5.16
\\
9
&
4.45
\\
10
&
3.98
\\
11
&
3.27
\\
12
&
3.07
\end{array}
\]

%%%%%%%%%%%%%%%%%%%%%%%%%%%%%%%%%%%%%%%%%%%%%%%%%%%%%%%%%%%%%%

\section*{Ein bemerkenswertes Resultat}

Besonders auffällig war

\[
t
=
10.
\]

Hier ergab sich

\[
d_{\mathrm{eff}}
=
3.98
\pm
0.36.
\]

Zum ersten Mal erschien die gesuchte Vier,

ohne dass sie explizit in den Algorithmus eingebaut worden war.

Dies war einer der wichtigsten numerischen Befunde der gesamten Untersuchung.

%%%%%%%%%%%%%%%%%%%%%%%%%%%%%%%%%%%%%%%%%%%%%%%%%%%%%%%%%%%%%%

\section*{Ist die Vier universell?}

Natürlich stellte sich sofort die nächste Frage.

War

\[
t
=
10
\]

eine reine Zufallszahl?

Oder existiert ein universelles Skalengesetz?

Falls Letzteres zutrifft,

müsste der optimale Wert von

\[
t
\]

mit der Größe des Systems skalieren.

%%%%%%%%%%%%%%%%%%%%%%%%%%%%%%%%%%%%%%%%%%%%%%%%%%%%%%%%%%%%%%

\section*{Der Skalierungsansatz}

Als erste Vermutung wurde angenommen,

dass die optimale Diffusionszeit proportional zum Logarithmus der Punktzahl wächst.

Es wurde daher untersucht:

\[
\boxed{
t
=
c
\,
\log N.
}
\]

Hierbei bezeichnet

\[
N
\]

die Zahl der Attraktoren,

während

\[
c
\]

eine dimensionslose Konstante ist.

%%%%%%%%%%%%%%%%%%%%%%%%%%%%%%%%%%%%%%%%%%%%%%%%%%%%%%%%%%%%%%

\section*{Der universelle Skalenscan}

Für verschiedene Werte von

\[
N
\]

und

\[
c
\]

wurden umfangreiche Simulationen durchgeführt.

Typische Resultate lauteten:

\[
\begin{array}{c|c|c}
N
&
c
&
d_{\mathrm{eff}}
\\
\hline
20
&
3.0
&
5.22
\\
40
&
2.5
&
4.50
\\
40
&
3.0
&
3.58
\\
80
&
2.0
&
3.86
\\
160
&
1.5
&
4.29
\end{array}
\]

%%%%%%%%%%%%%%%%%%%%%%%%%%%%%%%%%%%%%%%%%%%%%%%%%%%%%%%%%%%%%%

\section*{Interpretation}

Die Ergebnisse zeigten,

dass kein einzelner universeller Wert

\[
c
\]

existiert.

Gleichzeitig entstand jedoch ein anderer Eindruck.

Für geeignete Kombinationen aus

\[
N
\]

und

\[
t
\]

taucht immer wieder

\[
d_{\mathrm{eff}}
\approx
4
\]

auf.

Die Vier erscheint somit nicht an einer einzigen Stelle,

sondern auf einer ganzen Kurve im Parameterraum.

%%%%%%%%%%%%%%%%%%%%%%%%%%%%%%%%%%%%%%%%%%%%%%%%%%%%%%%%%%%%%%

\section*{Physikalische Deutung}

Dieses Verhalten besitzt eine bemerkenswerte Interpretation.

Der Parameter

\[
t
\]

ist keine fundamentale Naturkonstante.

Er beschreibt vielmehr,

wie weit ein Beobachter Informationen innerhalb des OPP integriert.

Kleine Werte entsprechen lokalen Beobachtern.

Große Werte beschreiben zunehmend globale Beobachtungen.

Die Raumzeit erscheint somit erst ab einer bestimmten Integrationsskala.

%%%%%%%%%%%%%%%%%%%%%%%%%%%%%%%%%%%%%%%%%%%%%%%%%%%%%%%%%%%%%%

\section*{Ein neuer Gedanke}

An dieser Stelle entstand erstmals eine weitreichende Hypothese.

Vielleicht ist

\[
t
\]

gar kein mathematischer Hilfsparameter.

Vielleicht beschreibt er unmittelbar die Auflösung des Beobachters.

Dann wäre

\[
\boxed{
t
=
\text{Informationsreichweite eines Beobachters}.
}
\]

Die Raumzeit wäre damit nicht nur projektional,

sondern zusätzlich skalenabhängig.

%%%%%%%%%%%%%%%%%%%%%%%%%%%%%%%%%%%%%%%%%%%%%%%%%%%%%%%%%%%%%%

\section*{Konsequenzen}

Die Untersuchungen führten zu mehreren wichtigen Schlussfolgerungen.

\begin{enumerate}

\item

Die Diffusionsprojektion besitzt tatsächlich einen physikalisch relevanten Parameter.

\item

Vierdimensionale Raumzeit erscheint nur innerhalb eines bestimmten Skalenbereiches.

\item

Die Raumzeit ist daher keine fundamentale Eigenschaft des OPP,

sondern eine emergente Beobachtungsstruktur.

\item

Der Beobachter wird damit zu einem unverzichtbaren Bestandteil der Theorie.

\end{enumerate}

%%%%%%%%%%%%%%%%%%%%%%%%%%%%%%%%%%%%%%%%%%%%%%%%%%%%%%%%%%%%%%

\section*{Zusammenfassung}

Die Diffusionsanalysen lieferten einen der stärksten numerischen Hinweise der gesamten PST.

Sie zeigten,

dass ein kontinuierlicher Übergang existiert

\[
\boxed{
\text{hochdimensionaler OPP}
\longrightarrow
\text{vierdimensionale Raumzeit}.
}
\]

Dieser Übergang wird nicht durch eine Änderung des OPP selbst erzeugt,

sondern ausschließlich durch die Art,

wie Informationen innerhalb des OPP integriert werden.

Damit erhält die Projektion erstmals eine direkte physikalische Bedeutung.

Sie beschreibt nicht nur einen mathematischen Algorithmus,

sondern möglicherweise den fundamentalen Prozess,

durch den Beobachter überhaupt eine Raumzeit erfahren.

